\section{Semidirect products and orbit intersections}\label{sec:semidirect}

Let $\alpha=(A,(A_i)_{i\in S})$ and
$\beta=(B,(B_j)_{j\in T})$ be flag-transitive coset geometries of
finite rank. Write $H_A=A_S$ and $H_B=B_T$. Suppose that $B$ acts on
$A$ by automorphisms and that an indexed action on $S$ is specified,
with $b(A_i)=A_{bi}$. This specification also covers coincident
parabolics: an action on an unlabelled set of subgroups alone does not
automatically specify an action on their repeated indices. Let
$\mathcal O$ be the orbit partition of $S$.

For an orbit $s\in\mathcal O$, put $A_s=\bigcap_{i\in s}A_i$.
In $G=A\rtimes B$ take the indexed parabolics
\[
 P_s=A_s\rtimes B\quad(s\in\mathcal O),\qquad
 P_j=A\rtimes B_j\quad(j\in T).
\]
The orbit types here use the whole orbit intersection. The construction below instead uses orbits of parabolic subgroups; these are different incidence structures with different panel formulas.

\begin{proposition}\label{J-prop:join}
Let $\overline\alpha=(A,(A_s)_{s\in\mathcal O})$. The coset
incidence system $(G,(P_i))$ is the join
$\overline\alpha*\beta$: the two component incidences are unchanged
and every element of one component is incident with every element of
the other. Both $\overline\alpha$ and the join are flag-transitive
geometries. For $J=J_O\sqcup J_T$, put
$E=\bigcup_{s\in J_O}s$. Then
\[
 G_J=A_E\rtimes B_{J_T},\qquad G_{\mathcal O\sqcup T}=H_A\rtimes H_B.
\]
\end{proposition}
\begin{proof}
Every $A_s$ is $B$-invariant. In semidirect coordinates,
\[
 (a,b)P_s=(aA_s)\times B,\qquad
 (a,b)P_j=A\times(bB_j).
\]
Here the second equality uses that $b$ acts bijectively on $A$.
Thus an orbit-type element is determined by the coset $aA_s$, a
$T$-type element by $bB_j$, and their incidences are exactly those
stated. The parabolic formula follows by intersecting the two
coordinates; this is ordinary intersection of Cartesian products.

The cosets of the $A_s$ belong to the family of cosets of standard
parabolics of $\alpha$. That family is Helly by Lemma~\ref{F-lem:cosets}. Hence $A$ is flag-transitive on $\overline\alpha$, and every
flag there extends to a chamber. The same is assumed for $\beta$.
For a flag of the join, choose $b\in B$ carrying its $\beta$ part
to the standard flag. This acts on its $\overline\alpha$ part;
then an element of $A$ carries that part to the standard flag while
fixing the $\beta$ part. Thus $G$ is flag-transitive. Alternatively,
a common $B$ coordinate and a common $A$ coordinate in the displayed
cosets give a common group representative directly.
\end{proof}

\begin{proposition}\label{J-prop:residues}
If both inputs are residually connected, so is the semidirect join.
Its panel indices are
\[
 [A_{S\setminus s}:H_A]\quad(s\in\mathcal O),\qquad
 [B_{T\setminus\{j\}}:H_B]\quad(j\in T).
\]
It is firm when both inputs are firm. If both inputs are thin, the
join is thin exactly when every orbit in $S$ is a singleton.
\end{proposition}
\begin{proof}
For clarity we first prove residual connectivity of
$\overline\alpha$. Put $R_i=A_{S\setminus\{i\}}$.
Residual connectivity of the finite-rank flag-transitive input gives
\[
 A_{S\setminus E}=\langle H_A,R_i:i\in E\rangle
 \qquad(E\subseteq S).
\]
This is Lemma~\ref{O-lem:parabolic-join}. Take a coarse
residue with at least two orbit types remaining, whose union is $E$.
For every $i\in E$ choose a remaining orbit $s$ not containing $i$.
Then $R_i\leq A_{(S\setminus E)\cup s}$, one of the maximal
parabolics of that residue. Those parabolics also contain $H_A$ and
therefore generate its whole group $A_{S\setminus E}$. This proves
residual connectivity of the coarse geometry.

A residue of a join is the join of the corresponding component
residues. If both components have a type remaining, the incidence
graph is connected because all vertices across the two components
are incident. If just one component has types remaining, residual
connectivity is that of the corresponding input residue. This proves
the first assertion, including empty component type sets.

The panel formulas follow by omitting just one type from the
parabolic formula in Proposition~\ref{J-prop:join}. For an orbit $s$,
choose $i\in s$. Firmness of $\alpha$ gives $H_A<R_i$, and
$R_i\leq A_{S\setminus s}$, so its panel index is at least two.
The other panels have exactly their $\beta$ sizes.

Suppose both inputs are thin. Singleton orbit indices are two. If
$s$ contains distinct $i,k$, choose $x\in R_i\setminus H_A$ and
$y\in R_k\setminus H_A$. We have $R_i\cap R_k=H_A$ by their
definitions. The four left cosets
\[
 H_A,\quad xH_A,\quad yH_A,\quad xyH_A
\]
are distinct. For example, $xyH_A=xH_A$ would put $y$ in $H_A$;
$xyH_A=yH_A$ would put $x\in yH_Ay^{-1}\leq R_k$, a
contradiction; and $xyH_A=H_A$ would put $x\in R_k$.
Thus $[A_{S\setminus s}:H_A]\geq4$, ruling out thinness.
The argument uses only subgroup intersections and does not assume
that $H_A$ is normal in all of $A$.
\end{proof}


\subsection{A general orbit construction}
Let $\alpha=(A,(A_x)_{x\in S})$ and
$\beta=(B,(B_i)_{i\in I})$ be flag-transitive coset geometries.
Suppose a specified homomorphism $B\to\Aut(A)$ and an action of $B$
on the indexed set $S$ satisfy $b(A_x)=A_{bx}$.
The indexed action is part of the data, even if different maximal
parabolics happen to coincide. Let $\mathcal K$ be the set of orbits
on $S$, choose $x_L\in L$ for each $L\in\mathcal K$, and assume
\begin{equation}\label{O-eq:general-orbits}
 B_Hx_L\cap B_Kx_L=B_{H\cup K}x_L
 \quad(H,K\subseteq I,\ L\in\mathcal K).
\end{equation}
For $T=I\sqcup\mathcal K$, define subgroups of $G=A\rtimes B$ by
\begin{equation}\label{O-eq:general-maximal}
 G_i=A_{S\setminus\bigcup_L B_ix_L}\rtimes B_i\quad(i\in I),
 \qquad G_L=A_L\rtimes B\quad(L\in\mathcal K).
\end{equation}
Every indicated subset of $S$ is invariant under its second factor,
so these are indeed subgroups. For $J\subseteq T$ write
$J_B=J\cap I$, $J_A=J\cap\mathcal K$, and set
\[
 E(J)=S\setminus\bigcup_{L\in\mathcal K\setminus J_A}B_{J_B}x_L.
\]

\begin{theorem}[Orbit construction with a Borel subgroup]\label{O-thm:general-twist}
The coset system in (\ref{O-eq:general-maximal}) is a flag-transitive
geometry. All its standard parabolics, including its Borel subgroup, are
\begin{equation}\label{O-eq:general-parabolics}
 G_J=A_{E(J)}\rtimes B_{J_B},\qquad G_T=A_S\rtimes B_I.
\end{equation}
It is finite if and only if both input geometries are finite. It is
residually connected if both inputs are residually connected, and firm
if both inputs are firm. Its numbers of chambers through panels are
\begin{align}
 [G_{T\setminus\{i\}}:G_T]&=[B_{I\setminus\{i\}}:B_I]
       &&(i\in I),\label{O-eq:old-panel}\\
 [G_{T\setminus\{L\}}:G_T]&=[A_{S\setminus B_Ix_L}:A_S]
       &&(L\in\mathcal K).\label{O-eq:new-panel}
\end{align}
In particular, if both inputs are thin and $B_I$ fixes each chosen
$x_L$, the new geometry is thin. Faithfulness of either input is not
needed for this sufficient condition.
\end{theorem}
\begin{proof}
Intersections of subgroups $A_E\rtimes B_H$ are computed coordinate by
coordinate. The intersection of their $A$-parts is obtained by taking
the union of their subscripts. Different $B$-orbits are disjoint, and
iteration of (\ref{O-eq:general-orbits}) says that the intersection of
the old supports is the support for $B_{J_B}$. Intersecting with the
orbit parabolics deletes precisely the orbits in $J_A$. Taking
complements proves (\ref{O-eq:general-parabolics}). It uses intersections
of maximal parabolics, not an unjustified formula for generated supports.

To prove flag transitivity, project a flag of $G$-cosets to $B$.
Old types give pairwise-intersecting maximal-parabolic cosets; orbit
types project to all of $B$. Coset intersection from the first section
gives a common $b\in B$. After left translation by $(1,b^{-1})$,
each coset has a representative $(a,1)$. Such a coset is the Cartesian
product of a coset of a standard parabolic of $A$ and the corresponding
subgroup of $B$. Pairwise incidence therefore gives a pairwise-intersecting
family of standard-parabolic $A$-cosets. Their common representative,
again by the coset intersection lemma, realizes the entire flag.
Thus every flag is an image of a standard flag and extends to a chamber.
This step by itself does not require the orbit identity; that identity
is needed for the exact parabolic description and the later conclusions.

Suppose both inputs are residually connected, and fix $J$ with at
least two missing types. We prove generation of $G_J$ by the next
parabolics. If no orbit type is missing, all their $A$-parts are $A_S$;
(\ref{O-eq:rc-borel}) for $B$ gives the required generation. If no old
type is missing, choose two distinct missing orbits $L,M$. The two
corresponding next parabolics both contain $B_I$. Their $A$-subscripts
are $E(J)\cup L$ and $E(J)\cup M$, whose intersection is $E(J)$.
Lemma~\ref{O-lem:parabolic-join} gives the entire $A$-part.

In the remaining case choose a missing old type $i$ and a missing
orbit type $L$. The parabolic $G_{J\cup\{L\}}$ contains the full
$B_{J_B}$. Write $E'=E(J\cup\{i\})$. The group generated by these
two next parabolics therefore contains $A_{bE'}$ for every
$b\in B_{J_B}$. In each remaining orbit, $B_{J_B\cup\{i\}}x_M$
is nonempty and contains $x_M$. Its $B_{J_B}$-saturation is exactly
$B_{J_B}x_M$. Consequently
\[
 \bigcap_{b\in B_{J_B}}bE'=E(J).
\]
There are only finitely many distinct subsets of the finite set $S$,
so Lemma~\ref{O-lem:parabolic-join} applies directly to these conjugate
parabolics. They generate $A_{E(J)}$, proving the remaining case.
The residue identification and connectivity criterion already proved
now establish residual connectedness in every residue.

Putting $J=T\setminus\{t\}$ in
(\ref{O-eq:general-parabolics}) gives the two panel formulas.
For an old type the $A$-parts agree; for an orbit type the $B$-parts
agree. These observations prove the stated indices even when the
Borel subgroups are nontrivial. If the inputs are firm, the old indices
are at least two. The orbit $B_Ix_L$ contains $x_L$, so
$A_{S\setminus\{x_L\}}\leq A_{S\setminus B_Ix_L}$; its index over
$A_S$ is at least two by firmness of $\alpha$. This proves firmness.
If the inputs are thin and $B_Ix_L=\{x_L\}$, all the displayed
indices are two, proving the sufficient thinness assertion.

Finally the number of chambers is the index of $G_T$. The elementary
coset index formula gives
\[
 [G:G_T]=[A:A_S][B:B_I],
\]
with the equality interpreted as a product of cardinalities.
For a finite number of types, a geometry has finitely many elements
if and only if it has finitely many chambers: every element extends
to a chamber and a chamber is a tuple of one element of each type.
The same argument applies to each input. Thus the product is finite
exactly when both input geometries are finite. No finiteness of the
possibly unfaithful groups $A$ and $B$ was assumed.
\end{proof}

\begin{corollary}[Faithful regular inputs]\label{O-thm:twist}
For faithful regular hypertopes $\alpha,\beta$, the construction is a
regular hypertope with full type-preserving automorphism group $A\rtimes B$.
Writing $A(R)=\langle a_x:x\in R\rangle$, its parabolics are
\begin{equation}\label{O-eq:regular-support}
 G_J=A\!\left(\bigcup_{L\notin J_A}B_{J_B}x_L\right)\rtimes B_{J_B}.
\end{equation}
Its distinguished involutions are the $b_i$ and the $a_{x_L}$.
\end{corollary}
\begin{proof}
The faithful regular model has $A_S=B_I=1$ and
$A_E=A(S\setminus E)$, as proved in the first section. The preceding
theorem supplies flag transitivity, residual connectedness and thinness.
Its Borel subgroup is trivial, and the panel groups are precisely the
cyclic groups on the indicated involutions. Chamber propagation shows
that the displayed regular action is the full type-preserving group.
\end{proof}
This recovers the general finite-type twisting theorem and its regular
specialization in Piedade--Tranchida~\cite[Theorem 3.16 and Corollary 3.17]{N006}.
The exact panel formula also explains why passing to a faithful quotient
of $B$ cannot silently preserve an arbitrary action on the types of
$\alpha$: its Borel subgroup must act trivially for that action to descend.
